% ID: SEL-D0036-OB-E01
% ESTADO: validada. S037, cotejada con examen y criterios.
\textbf{a)} El recinto queda entre la cota inferior
$\max(0,5/2-2x)$ y la superior $\min(2-2x/5,3-x)$.
En el eje horizontal se obtiene $5/4\le x\le3$.
La primera frontera inferior corta a $2-2x/5$ en
$(5/16,15/8)$. Las dos superiores se cortan en $(5/3,4/3)$.
En $x=0$ la cota inferior supera la superior, así que
no hay vértice sobre el eje vertical.
\begin{center}
\setlength{\unitlength}{12mm}
\begin{picture}(6,4)(-1,-1)
\qbezier(0,0)(1.8,0)(3.6,0)\put(3.7,0){$x$}
\qbezier(0,0)(0,1.15)(0,2.3)\put(0,2.4){$y$}
\qbezier(0.3125,1.875)(0.78125,0.9375)(1.25,0)
\qbezier(0.3125,1.875)(0.98958,1.60417)(1.66667,1.33333)
\qbezier(1.66667,1.33333)(2.33333,0.66667)(3,0)
\put(0.42,2.03){$(5/16,15/8)$}
\put(1.8,1.5){$(5/3,4/3)$}
\put(1.25,-0.15){\makebox(0,0)[t]{$(5/4,0)$}}
\put(3,-0.15){\makebox(0,0)[t]{$(3,0)$}}
\end{picture}
\end{center}
\textbf{b)} En los vértices $(5/4,0)$, $(3,0)$, $(5/3,4/3)$
y $(5/16,15/8)$, la función lineal vale $5/4$, $3$, $13/3$
y $65/16$, respectivamente. El mínimo es $\boxed{5/4}$ sólo
en $(5/4,0)$ y el máximo es $\boxed{13/3}$ sólo en $(5/3,4/3)$.
